% ============================================================
%  SEKCJA 2: MIĘDZYNARODOWE RAMY REGULACYJNE
% ============================================================

\section{Międzynarodowe ramy regulacyjne: co standardy rzeczywiście zakazują}

Ocena jakości i~autentyczności miodu opiera się na~sieci wzajemnie
powiązanych dokumentów normatywnych o~zasięgu zarówno globalnym jak i~regionalnym.
Zrozumienie ich zakresu i~ograniczeń jest kluczowe dla~oceny, czy obecność drożdży osmofilnych może stanowić samodzielną podstawę
dyskwalifikacji produktu z~obrotu handlowego. 

\subsection{Codex Alimentarius, CODEX STAN 12-1981}

Podstawowym dokumentem regulującym jakość miodu na~poziomie
międzynarodowym jest Codex Alimentarius Standard for Honey
(CODEX STAN 12-1981, ostatnia rewizja 2023), opracowany wspólnie
przez FAO i~WHO~\cite{codex_stan12}. Standard ten definiuje miód
jako naturalny produkt wytwarzany przez pszczoły  miodne
z~nektaru roślinnego, wydzielin żywych części roślin lub wydzielin
owadów ssących, który pszczoły zbierają, transformują, łącząc z~własnymi specyficznymi substancjami, odwadniają, magazynują
i~pozostawiają do~dojrzewania w~plastrach~\cite{codex_stan12}.

W zakresie parametrów fizykochemicznych CODEX STAN 12-1981 ustanawia
następujące limity ilościowe: zawartość wody $\leq 20{,}0\%$ (w~przypadku
miodu wrzosowego \textit{Calluna} $\leq 23{,}0\%$), zawartość
hydroksymetylofurfuralu (HMF) $\leq 40$~mg/kg (lub $\leq 80$~mg/kg
dla~miodów z~krajów o~klimacie tropikalnym), aktywność diastazy
$\geq 8$~jednostek Schade (lub $\geq 3$~j.~Schade dla miodów
o~naturalnie niskiej aktywności enzymatycznej, jak miód z~akacji
\textit{Robinia pseudoacacia}) oraz zawartość cukrów
redukujących~\cite{codex_stan12}. Standard nie~zawiera \textbf{żadnego
limitu liczby drożdży wyrażonego w~CFU/g}. Jedynym kryterium
mikrobiologicznym jest zakaz obrotu miodem \emph{,,który uległ
fermentacji lub jest w~trakcie fermentacji''} (\textit{,,that has
begun to ferment or effervesce''})~\cite{codex_stan12}. Jest to
rozróżnienie o~fundamentalnym znaczeniu prawnym: kryterium
dyskwalifikującym jest \emph{stan aktywnej fermentacji}, a~nie
\emph{obecność zdolnych do~fermentacji mikroorganizmów}.

\subsection{Dyrektywa UE 2001/110/WE i prawo krajowe}

Na~poziomie Unii Europejskiej podstawowym aktem prawnym regulującym obrót miodem jest Dyrektywa Rady 2001/110/WE z~dnia 20~grudnia 2001~roku w~sprawie miodu, implementowana w~polskim porządku prawnym rozporządzeniem
Ministra Rolnictwa i~Rozwoju Wsi~\cite{eudir2001110}. 
Dyrektywa ta, wzorowana na~standardzie Codex, nie~ustanawia limitu CFU/g dla drożdży ani grzybów w~miodzie. Parametry mikrobiologiczne miodu podlegają natomiast przepisom ogólnym dotyczącym higieny żywności, w~szczególności Rozporządzeniu (WE) nr~852/2004 Parlamentu Europejskiego i~Rady
w~sprawie higieny środków spożywczych~\cite{eudir2001110}.

W~praktyce kontrolnej laboratoriów inspekcji weterynaryjnej i~sanitarnej
funkcjonuje nieoficjalny próg higieniczny wynoszący $100$~CFU/g dla
łącznej liczby grzybów i~drożdży, stosowany jako ogólny wskaźnik
higieny procesu produkcji, jednak próg ten nie~posiada umocowania
w~przepisach szczegółowych dotyczących miodu i~nie~może stanowić
samodzielnej podstawy do~stwierdzenia niezgodności produktu
z~wymaganiami prawa żywnościowego w~rozumieniu art.~14
Rozporządzenia~(WE)~nr~178/2002~\cite{eudir2001110}. Dyrektywa
2001/110/WE dopuszcza ponadto możliwość filtrowania, podgrzewania
(do~określonych temperatur) i~mieszania partii miodu jako legalnych
zabiegów technologicznych, o~ile nie~prowadzą do~zmiany jego
podstawowego charakteru~\cite{eudir2001110}.


\subsection{Farmakopea Stanów Zjednoczonych (USP) i inne standardy krajowe}

Farmakopea Stanów Zjednoczonych (United States Pharmacopeia, USP)
klasyfikuje miód jako substancję pomocniczą o~zastosowaniu
farmaceutycznym i~spożywczym. Wymagania USP obejmują parametry
fizykochemiczne zbliżone do~Codex, jednak również nie~zawierają limitu
CFU/g dla drożdży osmofilnych~\cite{tsagkaris2021}. Analogiczne podejście
przyjmują standardy krajowe obowiązujące w~największych krajach
eksportujących miód, Chinach (GB/T~18796-2012), Argentynie
(Código Alimentario Argentino, CAA) i~Brazylii (IN MAPA nr~89/2000), które definiują jakość miodu przez parametry chemiczne i~aktywność enzymatyczną, nie~przez liczbę drożdży~\cite{tsagkaris2021}.

\subsection{Zestawienie porównawcze standardów}

Tabela~\ref{tab:standardy} prezentuje porównanie kluczowych parametrów jakości miodu w~obowiązujących standardach międzynarodowych i~regionalnych pod~kątem kryteriów mikrobiologicznych.

\begin{table}[ht]
\begin{adjustwidth}{-\extralength}{0cm}
%} % If the paper is ``preprints'', please uncomment this parenthesis.
%\isPreprints{\begin{tabularx}{\textwidth}{CCCC}}{% This command is only used for ``preprints''.
		
\centering
\caption{Porównanie parametrów jakości miodu w~głównych standardach
         międzynarodowych. Skróty: n.d., nie dotyczy (brak limitu
         w~dokumencie); CFU, jednostki tworzące kolonie;
         HMF, hydroksymetylofurfural.}
\label{tab:standardy}
\small
%\begin{tabular}{lcccccc}
\begin{tabularx}{\fulllength}{CCCCCC}
\hline
\textbf{Standard} &
\textbf{Wilgotność} &
\textbf{HMF} &
\textbf{Diastaza} &
\textbf{Limit CFU/g} &
\textbf{Kryterium} \\
 & (\%) & (mg/kg) & (j.~Schade) & \textbf{drożdże} & \textbf{mikrobiolog.} \\
\hline
Codex STAN 12-1981       & $\leq 20{,}0$ & $\leq 40$  & $\geq 8$   & \textbf{brak} & zakaz fermentacji \\
Dyrektywa UE 2001/110/WE & $\leq 20{,}0$ & $\leq 40$  & $\geq 8$   & \textbf{brak} & zakaz fermentacji \\
USP (honey monograph)    & $\leq 21{,}0$ & n.d.       & n.d.        & \textbf{brak} & zakaz fermentacji \\
Rozp. (WE) 852/2004      & n.d.          & n.d.       & n.d.        & $\leq 100$    & ogólna higiena$^{*}$ \\
\hline
\multicolumn{6}{l}{$^{*}$~Limit 100~CFU/g (grzyby + drożdże łącznie) ma~charakter higieny ogólnej} \\
\multicolumn{6}{l}{\phantom{$^{*}$~}i~nie~jest kryterium specyficznym dla~miodu.}
\end{tabularx}
\end{adjustwidth}
\end{table}

\subsection{Wniosek regulacyjny}

Analiza powyższych dokumentów prowadzi do~jednoznacznego wniosku:
żaden obowiązujący standard krajowy ani~międzynarodowy nie~traktuje
\emph{obecności drożdży osmofilnych per~se} jako podstawy
dyskwalifikacji miodu z~obrotu handlowego. Kryterium
dyskwalifikującym jest wyłącznie \emph{wynik biologiczny}
(aktywna fermentacja, eferwescencja), a~nie \emph{przyczyna potencjalna}
(obecność mikroorganizmów zdolnych do~fermentacji). Rozróżnienie
to~ma~fundamentalne znaczenie dla~oceny praktycznej: miód zawierający
drożdże osmofilne w~liczbie dowolnej, lecz niewykazujący oznak
aktywnej fermentacji i~spełniający kryterium wilgotności $\leq 20\%$,
spełnia wymagania wszystkich obowiązujących standardów~\cite{codex_stan12,eudir2001110}.
Sam pomiar CFU/g bez~jednoczesnej oceny aktywności wody ($a_w$),
temperatury i~czasu przechowywania jest metodologicznie
niewystarczający do~orzekania o~niezgodności produktu
z~wymaganiami prawa żywnościowego~\cite{tsagkaris2021,schoder2026}.
